\section{World--Track Coordinate Transformation}

% ============================================================
\subsection{Motivation}
% ============================================================

\begin{remark}
Railway alignments are modeled in a one-dimensional intrinsic coordinate
system defined by arc-length.
\end{remark}

\begin{remark}
Real-world data, however, is given in global coordinate systems such as
WGS84, UTM, or local engineering reference systems.
\end{remark}

\begin{remark}
Thus, a transformation between world coordinates and track coordinates
is required.
\end{remark}

% ============================================================
\subsection{Track Coordinate System}
% ============================================================

\begin{definition}[Track coordinate system]
Let \(\gamma(s)\) denote the alignment curve parameterized by arc-length.

The track coordinate system at position \(s\) is defined by:
\begin{itemize}
\item longitudinal coordinate: \(s\),
\item lateral offset: \(d\),
\item vertical offset: \(h\).
\end{itemize}
\end{definition}

\begin{remark}
In the planar case, only \((s,d)\) is considered.
\end{remark}

% ============================================================
\subsection{Forward Transformation (Track to World)}
% ============================================================

\begin{definition}[Track-to-world mapping]
Given \((s,d)\), the corresponding world position is:
\[
x(s,d)
=
\gamma(s)
+
d \cdot \mathbf{n}(s),
\]
where \(\mathbf{n}(s)\) is the normal vector of the alignment.
\end{definition}

\begin{remark}
The tangent vector is:
\[
\mathbf{t}(s)
=
(\cos\theta(s), \sin\theta(s)),
\]
and the normal vector:
\[
\mathbf{n}(s)
=
(-\sin\theta(s), \cos\theta(s)).
\]
\end{remark}

---

% ============================================================
\subsection{Inverse Transformation (World to Track)}
% ============================================================

\begin{definition}[World-to-track mapping]
Given a world point \(x\), the track coordinates \((s,d)\) are defined by:
\[
(s^\ast, d^\ast)
=
\arg\min_{s,d}
\| x - (\gamma(s) + d\,\mathbf{n}(s)) \|.
\]
\end{definition}

\begin{remark}
This is a nonlinear projection problem onto the alignment.
\end{remark}

\begin{remark}
In practice, it is solved by:
\begin{itemize}
\item nearest point search,
\item local Newton iteration,
\item segment-wise projection.
\end{itemize}
\end{remark}

---

% ============================================================
\subsection{Local Approximation}
% ============================================================

\begin{remark}
For small offsets, the projection can be approximated by:
\[
d \approx (x - \gamma(s)) \cdot \mathbf{n}(s).
\]
\end{remark}

\begin{remark}
The arc-length coordinate is approximated by projection onto the tangent:
\[
s \approx \arg\min_s \| x - \gamma(s) \|.
\]
\end{remark}

---

% ============================================================
\subsection{Relation to Curvature}
% ============================================================

\begin{remark}
The transformation depends on curvature:
\begin{itemize}
\item high curvature increases projection sensitivity,
\item low curvature yields stable projections.
\end{itemize}
\end{remark}

\begin{remark}
Errors in \(s\) propagate nonlinearly into world space.
\end{remark}

---

% ============================================================
\subsection{Computational Cost}
% ============================================================

\begin{remark}
The inverse mapping \(x \mapsto (s,d)\) is computationally expensive:
\begin{itemize}
\item nonlinear,
\item iterative,
\item dependent on curve representation.
\end{itemize}
\end{remark}

\begin{remark}
Thus, it must be used selectively in large-scale systems.
\end{remark}

---

% ============================================================
\subsection{Level-of-Detail Strategy}
% ============================================================

\begin{definition}[LOD-dependent projection]
The world-to-track transformation is evaluated depending on scale:
\begin{itemize}
\item coarse scale: approximate projection,
\item fine scale: exact projection.
\end{itemize}
\end{definition}

\begin{remark}
This reflects the fact that high precision is only required locally.
\end{remark}

---

% ============================================================
\subsection{Hybrid Strategy}
% ============================================================

\begin{remark}
A hybrid approach combines:
\begin{itemize}
\item precomputed samples for fast lookup,
\item local parametric refinement for accuracy.
\end{itemize}
\end{remark}

\begin{remark}
This mirrors the hybrid modelling approach used for curvature.
\end{remark}

---

% ============================================================
\subsection{Connection to CRS}
% ============================================================

\begin{remark}
The world coordinates depend on the chosen coordinate reference system
(CRS).
\end{remark}

\begin{remark}
Thus, the full transformation chain is:
\[
\text{CRS} \;\rightarrow\; \text{world} \;\rightarrow\; \text{track}.
\]
\end{remark}

---

% ============================================================
\subsection{Connection to Realization}
% ============================================================

\begin{remark}
The mapping depends on the realized alignment:
\[
\gamma_{\mathrm{real}}(s)
\neq
\gamma_{\mathrm{target}}(s).
\]
\end{remark}

\begin{remark}
Thus, world-to-track mapping should always refer to the realized
geometry.
\end{remark}

---

% ============================================================
\subsection{Key Insight}
% ============================================================

\begin{proposition}
The world-to-track transformation is the central interface between
geometric modelling and real-world data.
\end{proposition}

---

% ============================================================
\subsection{Connection to AIM}
% ============================================================

\begin{summary}
The world--track transformation links the intrinsic alignment model
to external coordinate systems and measured data.

Its nonlinear and scale-dependent nature makes it a critical component
for visualization, optimization, and data integration.

Efficient handling requires hybrid and LOD-based strategies, consistent
with the overall AIM framework.
\end{summary}
